\input{../tex-inputs/latex-leseansicht-vorspann}

\section[Gustav Schwarzkopf an Arthur Schnitzler, 6. 8. 1914]{L04248 Gustav Schwarzkopf an Arthur Schnitzler, 6. 8. 1914}
\nopagebreak\mylabel{L04248v}
\rehead{ }\normalsize\beginnumbering\briefempfaengerindex{Schnitzler, Arthur@\textsc{Schnitzler, Arthur}!zzzSchwarzkopf, Gustav@\emph{von Gustav Schwarzkopf}!1914-08-061@{6. 8. 1914}|(be}
\toendnotes[C]{\smallbreak\pagebreak[2]}
\correspDesc{Versand  durch Gustav Schwarzkopf am 6. 8. 1914 in Wien
\newline{}Übermittlung  am 6. 8. 1914 in Wien
\newline{}Umleitung  am 11. 8. 1914 in Celerina
\newline{}Erhalt  durch Arthur Schnitzler im Zeitraum [7. 8. 1914 – 11. 8. 1914?] in Bad Ragaz}\toendnotes[C]{\smallbreak}
\Standort{CUL, Schnitzler, B 96a.}
\physDesc{Bildpostkarte, 658 Zeichen
\newline{}Handschrift: schwarze Tinte, deutsche Kurrent
\newline{}Versand: 1) Stempel: »\nobreak{}\oindex{I., Innere Stadt@\textbf{I., Innere Stadt}, \emph{Verwaltungsgebiet}|pwk}1/1 Wien 8, 6. VIII. 14, 8\nobreak{}«.   2) Stempel: »\nobreak{}\oindex{Celerina@\textbf{Celerina}|pwk}Cel{[}erina{]} (Graubünden), 11. VIII. 14, XII\nobreak{}«.  3) mit schwarzer Tinte von unbekannter Hand mit schwarzer Tinte die drei Adresszeilen
                                 gestrichen und ersetzt durch »\noindent{}Hotel Lattmann\oindex{Hotel Lattmann@\textbf{Hotel Lattmann}, \emph{Hotel}|pw}{ / }\uline{Ragaz\oindex{Bad Ragaz@\textbf{Bad Ragaz}|pw}}.« 
\newline{}Schnitzler: mit Bleistift Vermerk: »\textsc{Schwarzk}{[}opf{]}« }\toendnotes[C]{\smallbreak}\pstart{}{\pb}I. Tiefer Graben 17\oindex{Wien@\textbf{Wien}!I., Innere Stadt@\textbf{I., Innere Stadt}!Tiefer Graben 17@\textbf{Tiefer Graben 17}, \emph{Wohngebäude}|pw}\pend{}\pstart{}II. 6.\pend{}{\bigskip}\pstart{}{\pb}Herrn\pend{}\pstart{}\textsc{D\textsuperscript{r} Arthur
                  Schnitzler}\pend{}\pstart{}\textsc{Celerina\oindex{Celerina@\textbf{Celerina}|pw}}\pend{}\pstart{}\textsc{(Schweiz
                     Engadin\oindex{Engadin@\textbf{Engadin}, \emph{Tal}|pw})}\pend{}\pstart{}\textsc{Cresta
                     Palace\oindex{Cresta Palace@\textbf{Cresta Palace}, \emph{Hotel}|pw}}\pend{}{\bigskip}
\pstart
           \noindent{}\centering{}{\pb}{[}Fotografie des Hauses Tiefer
                     Graben 17\oindex{Wien@\textbf{Wien}!I., Innere Stadt@\textbf{I., Innere Stadt}!Tiefer Graben 17@\textbf{Tiefer Graben 17}, \emph{Wohngebäude}|pw}. Eventuell blicken Gustav
                  und Max Schwarzkopf\pwindex{Schwarzkopf, Max 12.\,6.\,1857 Wien – 14.\,4.\,1928 ebd.@\textsc{Schwarzkopf, Max} (12.\,6.\,1857 Wien – 14.\,4.\,1928 ebd.), \emph{Rechtsanwalt}|pw} aus dem Fenster im 2. Stock.{]}\pend
           \vspace{1em}
\pstart
           \raggedleft{}{\pb}6. VIII. 14.\pend
           \vspace{0.5em}
\pstart
           Bin seit 30. hier. Leider große Hitze und Straßenlärm machen den Aufenthalt nicht
               angenehmer. Es{ }ſind aber viele hier, die eigentlich hier nichts zu{ }ſuchen haben, die
               nur die Nervoſität hergetrieben hat. Laſſen Sie{ }ſich nicht auch dazu verleiten. In
                  \textsc{Ischl\oindex{Bad Ischl@\textbf{Bad Ischl}|pw}} war ich auch. Einen ganzen Tag! Jetzt iſt mein Bruder\pwindex{Schwarzkopf, Max 12.\,6.\,1857 Wien – 14.\,4.\,1928 ebd.@\textsc{Schwarzkopf, Max} (12.\,6.\,1857 Wien – 14.\,4.\,1928 ebd.), \emph{Rechtsanwalt}|pwv} in \textsc{Weissenbach\oindex{Weißenbach am Attersee@\textbf{Weißenbach am Attersee}, \emph{Verwaltungsgebiet}|pw}}, er iſt aber faſt Alleinherrſcher im \textsc{Hôtel} – Wiſſen Sie{ }ſchon, daß auch \textsc{Hugo\pwindex{Hofmannsthal, Hugo von 1.\,2.\,1874 Wien – 15.\,7.\,1929 Rodaun@\textsc{Hofmannsthal, Hugo von} (1.\,2.\,1874 Wien – 15.\,7.\,1929 Rodaun), \emph{Schriftsteller}|pw}} einberufen iſt?\pend
           
\pstart
           Ihre Frau\pwindex{Schnitzler, Olga 17.\,1.\,1882 Wien – 13.\,1.\,1970 Lugano@\textsc{Schnitzler, Olga} (17.\,1.\,1882 Wien – 13.\,1.\,1970 Lugano), \emph{Schauspielerin, Sängerin}|pwv} und Sie herzlichſt grüßend,{\\[\baselineskip]}Ihr \spacefill\mbox{Gustav.}\pend
           \leftskip=0em{}\selectlanguage{ngerman}\endnumbering\briefempfaengerindex{Schnitzler, Arthur@\textsc{Schnitzler, Arthur}!zzzSchwarzkopf, Gustav@\emph{von Gustav Schwarzkopf}!1914-08-061@{6. 8. 1914}|)be}\mylabel{L04248h}
\begin{anhang}
\end{anhang}\newcommand{\dateiname}{L04248}\newcommand{\titel}{Gustav Schwarzkopf an Arthur Schnitzler, 6. 8. 1914}\newcommand{\editorInnen}{Herausgegeben von Jahnke, SelmaMüller, Martin Anton}\input{../tex-inputs/latex-leseansicht-abspann}
